\section{F. Low types, orientation, and a non-orientable test}

\begin{statement}{F.1}{Proved}{Low orientable types, peripheral structures, and centres}
For an orientable surface of genus $g$ with $n$ punctures and
$2g-2+n\leq0$, the possible types and discrete fundamental groups are

$$
(0,0):1,\qquad (0,1):1,\qquad (0,2):\mathbb Z,\qquad
 (1,0):\mathbb Z^2.
$$

Each group equals its centre. For $r=0,1,2$,

$$
\operatorname{Env}(\mathbb Z^r)_{\mathbb Q}
 =\mathbb G_a^r\times\widehat{\mathbb Z}^{\,r}_{\mathbb Q},
 \qquad
 \operatorname{Env}(\mathbb Z^r)_{\mathbb F_p}
 =\widehat{\mathbb Z}^{\,r}_{\mathbb F_p}.
$$

The finite lattice stage for $\mathbb Z^r$ is $\mathbb G_a^r$.
For a finite group $F$, $\operatorname{Env}(F)=\underline F$.
In particular the projective-plane group gives
$\operatorname{Env}(C_2)=\underline{C_2}$, not $\mu_2$.

The types $(0,3)$ and $(1,1)$ both have discrete group $F_2$ and
trivial centre, but have different peripheral markings. The first
has three peripheral conjugacy classes represented by
$x,y,(xy)^{-1}$; the second has one represented by $[x,y]$,
up to the orientation convention. A comparison which forgets
these markings does not distinguish the two types.
These statements qualify the low types; they do not establish a
general anabelian reconstruction theorem.
\end{statement}
\begin{proof}
The inequality gives exactly the four listed pairs. A punctured
sphere retracts to a bouquet of $n-1$ circles when $n>0$.
The sphere is simply connected, and the torus cell decomposition
gives $\langle a,b:[a,b]=1\rangle=\mathbb Z^2$.
The first centre assertions follow immediately.

For a lattice, apply E.12 to $n\mathbb Z^r$. Its components are
indexed by $(\mathbb Z/n)^r and have coordinates representing
$a+nT$. Over $\mathbb Q$, $X=a+nT$ gives compatible isomorphisms
with $\mathbb G_a^r\times(\mathbb Z/n)^r$; taking the inverse
limit proves the generic-fibre formula. In characteristic $p$,
refinement from $n$ to $pn$ makes each polynomial coordinate
constant on the finer components, giving the profinite constant
fibre as in E.14. For finite $F$, the zero subgroup is its only
torsion-free abelian subgroup, so the lattice refinement system
has that single object. The standard cell presentation of the
projective plane has relation $a^2=1$. Constant $C_2$ is etale
at two, whereas $\mu_2$ is not.

The thrice-punctured sphere has the presentation
$\langle x,y,z:xyz=1\rangle$ and hence free group on $x,y$.
Removing a disc from a torus removes its commutator relation
and leaves a boundary curve represented by $[x,y]$.
The centralizer of a free generator $x$ in $F(x,y)$ is
$\langle x\rangle$: a reduced word containing a $y$-syllable
cannot commute with $x$, by comparing reduced normal forms.
The analogous assertion for $y$ and
$\langle x\rangle\cap\langle y\rangle=1$ prove that the centre
is trivial. The peripheral markings follow from the respective
cell decompositions. Quotienting any of the abelian low-type
groups by its centre would instead discard the whole group.
\end{proof}

\begin{statement}{F.2}{Proved}{Orientation reversal and the full manuscript group}
Let $\Pi_0,S$ be as in E.18 and put $J=\operatorname{diag}(-1,1)$.
Then

$$
S^\pm:=\operatorname{GL}_2(\mathbb Z)/[\Pi_0,\Pi_0]
       =S\rtimes C_2
$$

as a discrete group, with involution induced by conjugation by $J$.
The full manuscript finite quotient has order $24$, whereas that
for $S$ has order $12$. For the induced action on envelopes,

$$
\operatorname{Env}(S^\pm)
   =\operatorname{Env}(S)\rtimes\underline{C_2}.
$$

The same assertion holds for the finite lattice stage.
This does not assert an integral splitting of the lattice extension
by its order-$24$ quotient.

On real matrix points orientation is controlled by
$\operatorname{sgn}(\det)$; positive determinant is different from
determinant one. The real sign map is not asserted to be a
morphism of group schemes.
\end{statement}
\begin{proof}
Direct multiplication gives $J\lambda_iJ^{-1}=\lambda_i^{-1}$
for $i=0,1$. Thus $[\Pi_0,\Pi_0]$ is preserved by $J$ and by
$\operatorname{SL}_2(\mathbb Z)$ and is normal in
$\operatorname{GL}_2(\mathbb Z)$. The determinant map on the
latter has image $\{\pm1\}$ and section $-1\mapsto J$.
Taking the quotient gives the split orientation extension.
Its index-two subgroup has finite lattice quotient of order
$12$, so the full quotient has order $24$.

The sublattices $n\mathbb Z^2$ are stable under the involution.
Apply E.12 at each such stage and take the inverse limit.
The constant $C_2$ factor and its section persist.
The orientation section need not preserve the chosen finite
subgroup used to split E.18, so this argument does not supply
a section of the order-$24$ lattice quotient.

For the standard alternating form $\omega$ on a real plane,
$\omega(Av,Aw)=\det(A)\omega(v,w)$ by direct expansion.
An antiholomorphic involution of a complex torus reverses
orientation, and therefore has determinant $-1$ on integral
first homology, by the intersection pairing. This explains
the need for an orientation-reversing operator in a
general-linear comparison; it does not identify a real
arithmetic fundamental group with this split modular quotient.
The orientation group is constant $C_2$ throughout.
Replacing it by $\mu_2$ at two would change the group scheme.
\end{proof}

\begin{statement}{F.3}{Proved}{A nonsplit integral model for the Klein bottle}
Let

$$
K=\langle a,b:aba^{-1}=b^{-1}\rangle,\qquad
 L=\langle a^2,b\rangle\simeq\mathbb Z^2.
$$

Then $K/L=C_2$, the centre is $\langle a^2\rangle$, and the
finite lattice model of E.12 is the disjoint union of two
affine planes with multiplication

$$
(x,y,\epsilon)(x',y',\delta)=
 (x+x'+\epsilon\delta,\ y+(-1)^\epsilon y',\
       \epsilon+\delta\pmod2),
 \qquad \epsilon,\delta\in\{0,1\}.
$$

This extension by $\mathbb G_a^2$ has no section of its constant
$C_2$ quotient over $\mathbb Z$ or $\mathbb F_2$.
It splits over $\mathbb Z[1/2]$.
Thus even this non-orientable Euler-characteristic-zero case
requires the cocycle construction rather than an assumption
of an integral splitting.
\end{statement}
\begin{proof}
Normal forms $a^m b^n$ follow by moving each $a$ past each $b$.
They are distinct in the plane action

$$
a(x,y)=(x+1/2,-y),\qquad b(x,y)=(x,y+1).
$$

Thus $L$ is free abelian of rank two and has index two.
Conjugation by $a$ fixes $a^2$ and inverts $b$.
Choosing $a$ as representative of the nontrivial coset gives
$c(1,1)=(1,0)$, and E.12 gives the displayed group law.
To commute with $b$, a normal form must have even $a$ exponent;
to commute also with $a$, its $b$ exponent must be zero.
This proves the centre assertion.

Every point in the nonidentity component has square
$(2x+1,0,0)$. Over $\mathbb Z$ this is never the identity;
over $\mathbb F_2$ it is always $(1,0,0)$.
A section would require an element of order two in that
component. Over $\mathbb Z[1/2]$, $(-1/2,0,1)$ supplies one.

On $\mathbb C/(\mathbb Z+i\mathbb Z)$ the involution
$z\mapsto\overline z+1/2$ is antiholomorphic and fixed-point-free:
its real displacement cannot be an integer. Its square is
translation by one, hence the identity on the torus.
The quotient is the Klein bottle, with the displayed plane
transformations as deck generators. This supplies a concrete
orientation-reversing example and its finite lattice model.
It does not establish comparisons for arbitrary real
arithmetic actions or all non-orientable hyperbolic surfaces.
\end{proof}
